% Pursuing Stacks (cotes 134-2 to 134-8): archive cote 134-5, site page 111 (PDF page 112), typescript page 298.
% typescript. Draft compared: ps5.tex:333-387.
% AI-transcribed by Claude Opus 5.5 (claude-opus-5-5, adaptive thinking, effort high), Message Batches API, batch ps2 (msgbatch_01NVwr1Yuo8qJf5rykeGBXPK), 2026-11-12; not yet checked by eye.
\scanpage{134-5/111}
\mspage{298}

as the formal ingredient of a definition of ``elementary I-homotopy'' bet-%
ween two maps in M, without any assumptions whatever on \xout{M} the category
A we start with. The standart case is the one when I = \hand{$\Delta$}$_1$ , the ``unit
interval'', but never mind. The definition works just as well, when \hand{$\Delta$}
is replaced by any (let's say small) category \hand{$\Delta$} whatever. On the other
hand, it is immediate that for a functor M $\to$ \over{\unsure{M}}{N} as above, induced by
a functor A $\longrightarrow$ B, \xout{twoxmapsxinxMxwhichxare} for two map\hand{s} in M, any elemen-%
tary homotopy between them gives rise to an elementary homotopy \over{\unsure{on}}{of} their
images in \hand{N} - which \xout{xmpix} is just the well known fact (in case \hand{$\Delta$}=\hand{$\Delta$} ,
I = \hand{$\Delta$}$_1$ ) that M $\to$ \hand{N} is compatible with simplicial homotopy notions.

In order to fit this into the general framework of section 51, \xout{the}
let's remark that if A is a full subcategory of a category A', then for
a pair of maps in M, the elementary $\underline{I}$-homotopies between these are the
same as when considering the given maps as maps in M' = \uline{Hom}(\hand{$\Delta$} ,A'), in
which M is embedded as a full subcategory. Now any (small, say) catego-%
ry A can be embedded canonically into A' = A\^{} as a full subcategory,
and any functor \xout{fx (1)xxAxxxxBxxembeds}
\[ f:\ A \to B \]
embeds in the corresponding functor
\[ f_{!} :\ A\hat{} \to B\hat{} \quad , \]
hence the functor
\[ \hand{$\varphi$}:\quad M \longrightarrow N \]
embeds in the corresponding functor
\[ \hand{$\varphi'$}:\quad M' \longrightarrow N' \qquad , \quad M' = \underline{Hom}(\hand{$\Delta^{\circ}$},A\hat{}\,), \ N' = \underline{Hom}(\hand{$\Delta^{\circ}$},B\hat{}\,) \ . \]
\xout{Zxexfunctor} As A\^{}, B\^{} satisfy the required exactness properties, it fol-%
lows that the $\underline{I}$-homotopy notions in M',N' can be interpreted in terms of
the notions of section 51, with respect to $\underline{I}_A$ and $\underline{I}_B$ , \xout{interpretedxnow}
defined now as \inserted{(componentwise)} constant objects of A\^{} and B\^{} respectively. Still, f$_{!}$
commutes to finite products only if f does, so we are still left with ex-%
plaining why M' $\to$ N' is well behaved with respect to $\underline{I}$ - homotopy
notions. Equivalently, we \inserted[typed]{need only} \xout{\unsure{havexto}} see this in the case of (2) M$\to$ N,
when A and B are supposed to have the required exactness properties to
allow for the interpretation given above of the $\underline{I}$-homotopy notions in
terms of the formalism of section 51, and when moreover f (as f$_{!}$ above)

\margin{111}
